%% ============================================================
%% shared/dr_strategies.tex
%% Disaster Recovery Strategies + High Availability Patterns
%% Classification: BOTH (SAA covers basics; SAP goes deeper on cost/trade-offs)
%% ============================================================

\servicesection{Disaster Recovery \& High Availability}{\bothtag}{%
  RTO/RPO trade-offs and DR strategy selection across the cost/complexity spectrum.}

\subsection{Key Reliability Metrics}

\begin{keyfacts}
\begin{itemize}
  \item \textbf{RTO (Recovery Time Objective)} --- maximum acceptable downtime after a failure. Lower RTO = faster recovery = higher cost.
  \item \textbf{RPO (Recovery Point Objective)} --- maximum acceptable data loss, measured in time. Lower RPO = more frequent replication = higher cost.
  \item These metrics drive architecture decisions: always ask the business for RTO/RPO \emph{before} choosing a DR strategy.
\end{itemize}
\end{keyfacts}

% ─── TikZ: DR Strategy Spectrum ───────────────────────────────
\subsection{Disaster Recovery Strategies}

\begin{figure}[h]
\centering
\begin{tikzpicture}[
  box/.style={draw, rounded corners=4pt, minimum width=2.8cm, minimum height=2cm,
               align=center, font=\small},
  arrow/.style={->, thick, color=awsdark},
  label/.style={font=\footnotesize\itshape, color=gray}
]

%% Four DR strategy boxes arranged left (cheap/slow) to right (costly/fast)
\node[box, fill=saacolor!15,  draw=saacolor]  (br) at (0,0)
  {\textbf{Backup \&}\\\textbf{Restore}\\[4pt]\scriptsize Hours / Hours};

\node[box, fill=sharedcolor!15, draw=sharedcolor] (pl) at (3.4,0)
  {\textbf{Pilot}\\\textbf{Light}\\[4pt]\scriptsize Tens of mins\\/ Minutes};

\node[box, fill=tipcolor!20, draw=tipcolor] (ws) at (6.8,0)
  {\textbf{Warm}\\\textbf{Standby}\\[4pt]\scriptsize Minutes\\/ Seconds};

\node[box, fill=sapcolor!20, draw=sapcolor] (ms) at (10.2,0)
  {\textbf{Multi-Site}\\\textbf{Active/Active}\\[4pt]\scriptsize Near zero\\/ Near zero};

%% Arrows between boxes
\draw[arrow, color=sharedcolor] (br.east) -- (pl.west);
\draw[arrow, color=sharedcolor] (pl.east) -- (ws.west);
\draw[arrow, color=sharedcolor] (ws.east) -- (ms.west);

%% Cost axis (below)
\draw[->, thick, color=warncolor] (-1.0,-1.4) -- (11.4,-1.4)
  node[right, font=\small\bfseries, color=warncolor] {Cost \& Complexity};
\node[label] at (-1.0,-1.7) {Cheapest};
\node[label] at (11.2,-1.7) {Most Expensive};

%% RTO/RPO axis (above)
\draw[->, thick, color=saacolor, ->] (11.4,1.5) -- (-1.0,1.5)
  node[left, font=\small\bfseries, color=saacolor] {RTO / RPO};
\node[label] at (11.2,1.8) {Near Zero};
\node[label] at (-0.8,1.8) {Hours};

%% Labels
\node[font=\tiny, color=saacolor, below=2pt] at (br.south) {S3 / Glacier backup};
\node[font=\tiny, color=sharedcolor, below=2pt] at (pl.south) {Core services only};
\node[font=\tiny, color=tipcolor, below=2pt] at (ws.south) {Scaled-down live env};
\node[font=\tiny, color=sapcolor, below=2pt] at (ms.south) {Full prod in 2+ regions};

\end{tikzpicture}
\caption{DR Strategy Spectrum — RTO/RPO in each box; cost increases left to right.}
\end{figure}

\begin{comparebox}[title={DR Strategies}]
\begin{tabular}{@{}p{3.2cm}p{4cm}p{1.8cm}p{1.8cm}@{}}
\toprule
\textbf{Strategy} & \textbf{Description} & \textbf{RTO} & \textbf{RPO} \\
\midrule
Backup \& Restore    & Data backed up to S3/Glacier; restore from scratch    & Hours         & Hours        \\
Pilot Light          & Minimal core (DB replicated); scale up on failover     & Tens of mins  & Minutes      \\
Warm Standby         & Scaled-down but fully functional environment           & Minutes       & Seconds      \\
Multi-Site Active/Active & Full production in multiple regions simultaneously & Near zero     & Near zero    \\
\bottomrule
\end{tabular}
\end{comparebox}

\begin{examtip}
\textbf{AWS Elastic Disaster Recovery (DRS)} automates continuous block-level replication from on-prem or AWS to a target region. Use it for low-RTO/RPO without running full Warm Standby infrastructure.
Backup \& Restore is cheapest. Multi-Site is most expensive. \emph{The exam gives you RTO/RPO requirements — map them to a strategy.}
\end{examtip}

\begin{saadepth}
For SAA, focus on:
\begin{itemize}
  \item RDS Multi-AZ for intra-region HA (synchronous standby, auto failover, same region)
  \item Cross-Region Read Replicas — promote to standalone for DR in another region (async)
  \item S3 Cross-Region Replication (CRR) for object-level DR
  \item Route 53 health checks for DNS-based failover
\end{itemize}
\end{saadepth}

\begin{sapdepth}
At SAP level, expect to choose the \emph{right} strategy for given RTO/RPO/cost constraints:
\begin{itemize}
  \item \awsservice{Aurora Global Database} --- RPO $<$ 1 s; RTO $<$ 1 min. Primary region serves writes; up to 5 secondary regions for reads and DR promotion.
  \item \awsservice{DynamoDB Global Tables} --- multi-region, multi-active; built-in replication; sub-second RPO.
  \item \awsservice{AWS Backup} --- centralised backup across services (EC2, EBS, RDS, DynamoDB, EFS, FSx) with cross-account and cross-region copy.
  \item Pilot Light pattern: keep a minimal ``always-on'' config (IAM roles, DNS records, DB replica running); on disaster, scale up compute rapidly using pre-baked AMIs or CloudFormation.
  \item \awsservice{AWS Resilience Hub} --- continuously validates RTO/RPO against application architecture; recommends improvements.
\end{itemize}
\end{sapdepth}

% ─────────────────────────────────────────
\subsection{High Availability Patterns}

\begin{keyfacts}
\begin{itemize}
  \item \textbf{Multi-AZ deployments} --- protects against AZ failure (RDS Multi-AZ, ALB across AZs, Auto Scaling across AZs)
  \item \textbf{Auto Scaling Groups} --- automatically replace unhealthy instances; scale with demand; always span 2+ AZs
  \item \textbf{Elastic Load Balancing} --- distributes traffic; health checks; cross-zone load balancing (ALB default ON; NLB default OFF)
  \item \textbf{Route 53 health checks} --- DNS failover to healthy endpoints; can trigger CloudWatch alarm on failure
  \item \textbf{SQS decoupling} --- decouple components so failure of one does not cascade (producer keeps writing; consumer retries)
\end{itemize}
\end{keyfacts}

\begin{gotcha}
\textbf{ALB vs NLB cross-zone:} ALB distributes evenly across all targets in all enabled AZs by default (cross-zone ON). NLB and GLB have cross-zone load balancing \emph{off} by default --- enable it explicitly if you have uneven instance counts per AZ.
\end{gotcha}
